% ===== PRÉAMBULE CANONIQUE QCM — PHYSIQUE =====
% Sert à compiler controle.tex ET à la revalidation automatique du JSON
% (valider_qcm.py). NE PAS MODIFIER : packages figés.
% La bibliothèque TikZ "babel" est indispensable (conflit babel-french/circuitikz).
\documentclass[11pt,a4paper]{article}
\usepackage[utf8]{inputenc}\usepackage[T1]{fontenc}\usepackage{lmodern}
\usepackage[french]{babel}
\usepackage{amsmath,amssymb,mathtools}\usepackage{icomma}
\usepackage{siunitx}\sisetup{locale=FR}
\usepackage{xcolor}\usepackage{colortbl}
\usepackage{tikz}\usepackage{pgfplots}\pgfplotsset{compat=1.18}
\usepackage[siunitx,europeanresistors]{circuitikz}
\usepackage{enumitem}\usepackage{array,booktabs}
\usepackage[a4paper,margin=2.2cm]{geometry}
\usetikzlibrary{arrows.meta,calc,angles,quotes,positioning,patterns,%
  backgrounds,shapes.geometric,decorations.pathmorphing,%
  decorations.markings,intersections,babel}
\newcommand{\vv}[1]{\vec{#1}}\newcommand{\ds}{\displaystyle}
\newcommand{\R}{\mathbb{R}}\newcommand{\N}{\mathbb{N}}\newcommand{\Z}{\mathbb{Z}}
% ==============================================
\begin{document}
\begin{center}\Large\bfseries QCM concours --- Physique --- Fiche 09 : Électrostatique\end{center}
\vspace{6pt}\hrule\vspace{8pt}

% ================= Q01 =================
\noindent\textbf{PHYS-F09-Q01 --- Loi de Coulomb : force résultante nulle}\par
Deux charges ponctuelles fixes $q_A=+\qty{1}{\micro\coulomb}$ (en A) et $q_B=+\qty{9}{\micro\coulomb}$ (en B) sont séparées de $\qty{40}{\centi\metre}$. En quel point P de la droite $(AB)$ une charge test serait-elle en équilibre (force électrostatique résultante nulle) ?
\begin{center}
\begin{tikzpicture}[>=Stealth]
  \draw[thick] (0,0)--(8,0);
  \filldraw[red] (0,0) circle (2.5pt);
  \node[above=2pt] at (0,0){$q_A=+\qty{1}{\micro\coulomb}$};
  \node[below=3pt] at (0,0){A};
  \filldraw[red] (8,0) circle (2.5pt);
  \node[above=2pt] at (8,0){$q_B=+\qty{9}{\micro\coulomb}$};
  \node[below=3pt] at (8,0){B};
  \filldraw[black] (2.4,0) circle (2pt);
  \node[above=2pt] at (2.4,0){P\,?};
  \draw[<->] (0,-0.9)--(8,-0.9) node[midway,fill=white]{\qty{40}{\centi\metre}};
\end{tikzpicture}
\end{center}
\begin{enumerate}[label=\Alph*)]
\item à \qty{4}{\centi\metre} de A
\item à \qty{10}{\centi\metre} de A
\item à \qty{13}{\centi\metre} de A
\item à \qty{30}{\centi\metre} de A
\end{enumerate}
\textit{Bonne réponse : B}\par
Les deux charges étant de même signe, le point d'équilibre est entre elles, plus près de la plus faible ($q_A$). En notant $x$ la distance à A~: $\dfrac{Kq_A}{x^{2}}=\dfrac{Kq_B}{(d-x)^{2}}$, soit $\dfrac{\sqrt{q_A}}{x}=\dfrac{\sqrt{q_B}}{d-x}$, d'où $\dfrac{1}{x}=\dfrac{3}{40-x}$ et $x=\qty{10}{\centi\metre}$ (donc à \qty{30}{\centi\metre} de B).\par
\textit{Pièges :}
\begin{itemize}
\item A: équilibre écrit avec $q/d$ au lieu de $q/d^{2}$ (racine oubliée) : $1/x=9/(40-x)\Rightarrow x=\qty{4}{\centi\metre}$.
\item C: rapport $\sqrt{q_A/q_B}$ appliqué à $d$ au lieu de $(d-x)$ : $x=40/3\approx\qty{13}{\centi\metre}$.
\item D: rapport inversé, point rapproché de la grosse charge : $x=40\cdot 3/4=\qty{30}{\centi\metre}$.
\end{itemize}
\vspace{8pt}\hrule\vspace{8pt}

% ================= Q02 =================
\noindent\textbf{PHYS-F09-Q02 --- Loi de Coulomb : influence des charges et de la distance}\par
Deux charges ponctuelles exercent l'une sur l'autre une force d'intensité $F_0$. On triple \emph{simultanément} la valeur de chacune des deux charges \emph{et} la distance qui les sépare. La nouvelle intensité de la force vaut :
\begin{enumerate}[label=\Alph*)]
\item $F_0$ (inchangée)
\item $9\,F_0$
\item $F_0/9$
\item $3\,F_0$
\end{enumerate}
\textit{Bonne réponse : A}\par
$F=\dfrac{K q_1 q_2}{d^{2}}$. Tripler chaque charge multiplie le numérateur par $3\times 3=9$~; tripler la distance divise par $3^{2}=9$. Bilan : $F=\dfrac{9}{9}\,F_0=F_0$, l'intensité est \textbf{inchangée}.\par
\textit{Pièges :}
\begin{itemize}
\item B: effet des charges seul pris en compte ($\times 9$), variation de distance oubliée.
\item C: effet de la distance seul ($\div 9$), variation des charges oubliée.
\item D: distance traitée linéairement ($\div 3$ au lieu de $\div 9$) : $9/3=3$.
\end{itemize}
\vspace{8pt}\hrule\vspace{8pt}

% ================= Q03 =================
\noindent\textbf{PHYS-F09-Q03 --- Loi de Coulomb : direction de la résultante (carré)}\par
Trois charges identiques $+q$ occupent les sommets A, B et C d'un carré ; une quatrième charge $+q$ est placée au sommet D (charge étudiée). Dans quelle direction pointe la force électrostatique résultante subie par la charge en D ?
\begin{center}
\begin{tikzpicture}[>=Stealth,scale=1.4]
  \coordinate (A) at (0,2); \coordinate (B) at (2,2);
  \coordinate (C) at (2,0); \coordinate (D) at (0,0);
  \draw[gray,dashed] (A)--(B)--(C)--(D)--cycle;
  \draw[gray,dashed] (B)--(D);
  \filldraw[red] (A) circle (2.5pt) node[above left]{$+q$ (A)};
  \filldraw[red] (B) circle (2.5pt) node[above right]{$+q$ (B)};
  \filldraw[red] (C) circle (2.5pt) node[below right]{$+q$ (C)};
  \filldraw[teal] (D) circle (2.5pt) node[below left]{$+q$ (D)};
\end{tikzpicture}
\end{center}
\begin{enumerate}[label=\Alph*)]
\item verticalement, vers le bas
\item horizontalement, vers la gauche
\item selon la diagonale $(BD)$, vers l'extérieur du carré
\item selon la diagonale $(BD)$, vers le centre du carré
\end{enumerate}
\textit{Bonne réponse : C}\par
Les trois charges sont positives : chacune \emph{repousse} la charge en D. La charge A pousse D vers le bas, la charge C pousse D vers la gauche (mêmes intensités, distances égales à un côté), la charge B (diagonale) pousse D vers le bas-gauche. Par symétrie, la résultante est portée par la diagonale $(BD)$ et dirigée vers l'extérieur du carré (elle s'éloigne du centre).\par
\textit{Pièges :}
\begin{itemize}
\item A: seule la charge A (au-dessus) prise en compte : force verticale vers le bas.
\item B: seule la charge C (à droite) prise en compte : force horizontale vers la gauche.
\item D: sens inversé (attraction supposée) : diagonale dirigée vers le centre.
\end{itemize}
\vspace{8pt}\hrule\vspace{8pt}

% ================= Q04 =================
\noindent\textbf{PHYS-F09-Q04 --- Loi de Coulomb : superposition (triangle équilatéral)}\par
Trois charges identiques $+q$ sont placées aux sommets d'un triangle équilatéral de côté $a$. L'intensité de la force électrostatique résultante subie par l'une quelconque de ces charges vaut :
\begin{center}
\begin{tikzpicture}[scale=1.5]
  \coordinate (S) at (90:1.15);
  \coordinate (L) at (210:1.15);
  \coordinate (Rr) at (330:1.15);
  \draw[gray] (S)--(L)--(Rr)--cycle;
  \filldraw[red] (S) circle (2pt) node[above]{$+q$};
  \filldraw[red] (L) circle (2pt) node[below left]{$+q$};
  \filldraw[red] (Rr) circle (2pt) node[below right]{$+q$};
  \node at (0,-1.35){côté $a$};
\end{tikzpicture}
\end{center}
\begin{enumerate}[label=\Alph*)]
\item $F=2\,\dfrac{Kq^{2}}{a^{2}}$
\item $F=\dfrac{Kq^{2}}{a^{2}}$
\item $F=\sqrt{2}\,\dfrac{Kq^{2}}{a^{2}}$
\item $F=\sqrt{3}\,\dfrac{Kq^{2}}{a^{2}}$
\end{enumerate}
\textit{Bonne réponse : D}\par
Chaque charge subit deux forces de même intensité $F_1=\dfrac{Kq^{2}}{a^{2}}$, faisant entre elles l'angle du sommet, soit \ang{60}. La résultante vaut $F=2F_1\cos\!\left(\ang{30}\right)=2F_1\cdot\dfrac{\sqrt{3}}{2}=\sqrt{3}\,\dfrac{Kq^{2}}{a^{2}}$.\par
\textit{Pièges :}
\begin{itemize}
\item A: les deux forces sommées comme si elles étaient colinéaires ($2F_1$).
\item B: une seule interaction prise en compte, superposition oubliée ($F_1$).
\item C: angle traité comme droit (\ang{90}), théorème de Pythagore : $\sqrt{2}\,F_1$.
\end{itemize}
\vspace{8pt}\hrule\vspace{8pt}

% ================= Q05 =================
\noindent\textbf{PHYS-F09-Q05 --- Champ électrique : relation $\vv F=q\vv E$ (charge négative)}\par
Dans une région où règne un champ électrique uniforme $\vv E$ d'intensité $E=\qty{5.0}{\kilo\volt\per\metre}$ orienté vers la droite, on abandonne un électron (charge $-e$, avec $e=\qty{1.6e-19}{\coulomb}$). La force électrique subie par l'électron a pour intensité et sens :
\begin{enumerate}[label=\Alph*)]
\item $\qty{8.0e-16}{\newton}$, dans le sens opposé à $\vv E$ (vers la gauche)
\item $\qty{8.0e-16}{\newton}$, dans le sens de $\vv E$ (vers la droite)
\item $\qty{8.0e-19}{\newton}$, dans le sens opposé à $\vv E$
\item $\qty{1.6e-19}{\newton}$, dans le sens opposé à $\vv E$
\end{enumerate}
\textit{Bonne réponse : A}\par
$\vv F=q\vv E$ avec $q=-e$. Intensité : $F=eE=\num{1.6e-19}\times\num{5.0e3}=\qty{8.0e-16}{\newton}$. La charge étant négative, $\vv F$ est de sens \emph{opposé} à $\vv E$, donc vers la gauche.\par
\textit{Pièges :}
\begin{itemize}
\item B: intensité correcte mais signe de la charge oublié (force dans le sens de $\vv E$).
\item C: préfixe kilo négligé ($E=\qty{5.0}{\volt\per\metre}$) : $F=\qty{8.0e-19}{\newton}$.
\item D: champ oublié, la charge elle-même reportée comme une force : $\qty{1.6e-19}{\newton}$.
\end{itemize}
\vspace{8pt}\hrule\vspace{8pt}

% ================= Q06 =================
\noindent\textbf{PHYS-F09-Q06 --- Champ électrique créé par une charge ponctuelle}\par
Une charge ponctuelle $Q=+\qty{5.0}{\micro\coulomb}$ crée, en un point M situé à $r=\qty{0.30}{\metre}$, un champ électrique dont l'intensité vaut ($K=\qty{9.0e9}{\newton\metre\squared\per\coulomb\squared}$) :
\begin{enumerate}[label=\Alph*)]
\item $\qty{5.0e1}{\newton\per\coulomb}$
\item $\qty{4.5e4}{\newton\per\coulomb}$
\item $\qty{1.5e5}{\newton\per\coulomb}$
\item $\qty{5.0e5}{\newton\per\coulomb}$
\end{enumerate}
\textit{Bonne réponse : D}\par
$E=\dfrac{KQ}{r^{2}}=\dfrac{\num{9.0e9}\times\num{5.0e-6}}{(0{,}30)^{2}}=\dfrac{\num{4.5e4}}{0{,}09}=\qty{5.0e5}{\newton\per\coulomb}$.\par
\textit{Pièges :}
\begin{itemize}
\item A: distance laissée en centimètres et élevée au carré ($30^{2}=900$) : $\qty{5.0e1}{\newton\per\coulomb}$.
\item B: division par la distance oubliée (numérateur $KQ$ seul) : $\qty{4.5e4}{\newton\per\coulomb}$.
\item C: division par $r$ au lieu de $r^{2}$ : $\qty{4.5e4}{}/0{,}30=\qty{1.5e5}{\newton\per\coulomb}$.
\end{itemize}
\vspace{8pt}\hrule\vspace{8pt}

% ================= Q07 =================
\noindent\textbf{PHYS-F09-Q07 --- Condensateur plan : champ uniforme $E=U/d$}\par
Un condensateur plan est formé de deux disques de rayon $R=\qty{5.0}{\centi\metre}$ séparés d'une distance $d=\qty{4.0}{\milli\metre}$. Une tension $U=\qty{200}{\volt}$ est appliquée entre les armatures. L'intensité du champ électrique uniforme entre les armatures vaut :
\begin{center}
\begin{tikzpicture}[>=Stealth]
  \draw[thick,red] (0,1.4)--(3.2,1.4);
  \node[red,left] at (0,1.4){$+$};
  \draw[thick,blue] (0,0)--(3.2,0);
  \node[blue,left] at (0,0){$-$};
  \foreach \x in {0.5,1.1,1.7,2.3,2.9} \draw[->,teal] (\x,1.3)--(\x,0.1);
  \draw[<->] (3.6,0)--(3.6,1.4) node[midway,right]{$d$};
  \node at (1.6,1.85){rayon $R=\qty{5.0}{\centi\metre}$};
  \node at (1.6,-0.5){$U=\qty{200}{\volt}$};
\end{tikzpicture}
\end{center}
\begin{enumerate}[label=\Alph*)]
\item $\qty{5.0e1}{\volt\per\metre}$
\item $\qty{4.0e3}{\volt\per\metre}$
\item $\qty{5.0e4}{\volt\per\metre}$
\item $\qty{8.0e-1}{\volt\per\metre}$
\end{enumerate}
\textit{Bonne réponse : C}\par
Entre les armatures le champ est uniforme et vaut $E=\dfrac{U}{d}=\dfrac{200}{\num{4.0e-3}}=\qty{5.0e4}{\volt\per\metre}$. Le rayon $R$ des disques n'intervient pas.\par
\textit{Pièges :}
\begin{itemize}
\item A: distance $d$ laissée en millimètres ($E=200/4$) : $\qty{5.0e1}{\volt\per\metre}$.
\item B: rayon $R=\qty{5.0}{\centi\metre}$ utilisé comme distance : $200/0{,}05=\qty{4.0e3}{\volt\per\metre}$.
\item D: produit $U\times d$ au lieu du quotient : $200\times\num{4.0e-3}=\qty{8.0e-1}{\volt\per\metre}$.
\end{itemize}
\vspace{8pt}\hrule\vspace{8pt}

% ================= Q08 =================
\noindent\textbf{PHYS-F09-Q08 --- Travail de la force électrique (champ uniforme)}\par
Dans un champ électrique uniforme $E=\qty{2.0e3}{\volt\per\metre}$ orienté selon $+x$, une charge $q=+\qty{5.0}{\micro\coulomb}$ est déplacée de $d=\qty{0.20}{\metre}$ dans le sens $-x$ (c'est-à-dire contre le champ). Le travail de la force électrique lors de ce déplacement vaut :
\begin{enumerate}[label=\Alph*)]
\item $+\qty{2.0e-3}{\joule}$
\item $-\qty{2.0e-3}{\joule}$
\item $-\qty{1.0e-2}{\joule}$
\item $-\qty{2.0e-1}{\joule}$
\end{enumerate}
\textit{Bonne réponse : B}\par
La force est $F=qE=\num{5.0e-6}\times\num{2.0e3}=\qty{1.0e-2}{\newton}$, orientée selon $+x$ ($q>0$). Le déplacement se fait selon $-x$ : $W=F\,d\cos\ang{180}=-\num{1.0e-2}\times 0{,}20=-\qty{2.0e-3}{\joule}$ (travail résistant, négatif).\par
\textit{Pièges :}
\begin{itemize}
\item A: signe oublié, déplacement supposé dans le sens de la force ($\cos\ang{0}$).
\item C: multiplication par $d$ oubliée, la force reportée comme un travail : $-\qty{1.0e-2}{\joule}$.
\item D: distance laissée en centimètres ($d=20$) : $-\num{1.0e-2}\times 20=-\qty{2.0e-1}{\joule}$.
\end{itemize}
\vspace{8pt}\hrule\vspace{8pt}

% ================= Q09 =================
\noindent\textbf{PHYS-F09-Q09 --- Équilibre d'une gouttelette chargée (lévitation)}\par
Une gouttelette de masse $m=\qty{3.2e-15}{\kilogram}$ reste immobile entre deux plaques horizontales distantes de $d=\qty{2.0}{\centi\metre}$, soumises à une tension $U=\qty{400}{\volt}$. On prend $g=\qty{10}{\metre\per\second\squared}$. La valeur absolue de la charge portée par la gouttelette vaut :
\begin{center}
\begin{tikzpicture}[>=Stealth]
  \draw[thick,red] (0,1.6)--(3.2,1.6);
  \draw[thick,blue] (0,0)--(3.2,0);
  \filldraw[gray] (1.6,0.8) circle (3pt);
  \draw[->,teal,thick] (1.6,0.8)--(1.6,1.45) node[right]{$\vv F$};
  \draw[->,orange,thick] (1.6,0.8)--(1.6,0.15) node[right]{$\vv P$};
  \node[left] at (1.5,0.8){$m,\,q$};
  \draw[<->] (3.6,0)--(3.6,1.6) node[midway,right]{$d$};
\end{tikzpicture}
\end{center}
\begin{enumerate}[label=\Alph*)]
\item $\qty{1.6e-18}{\coulomb}$
\item $\qty{1.6e-16}{\coulomb}$
\item $\qty{1.6e-19}{\coulomb}$
\item $\qty{4.0e-15}{\coulomb}$
\end{enumerate}
\textit{Bonne réponse : A}\par
À l'équilibre la force électrique compense le poids : $|q|E=mg$ avec $E=\dfrac{U}{d}=\dfrac{400}{0{,}020}=\qty{2.0e4}{\volt\per\metre}$. D'où $|q|=\dfrac{mg}{E}=\dfrac{\num{3.2e-15}\times 10}{\num{2.0e4}}=\qty{1.6e-18}{\coulomb}$.\par
\textit{Pièges :}
\begin{itemize}
\item B: distance laissée en centimètres ($E=400/2=\qty{200}{\volt\per\metre}$) : $|q|=\qty{1.6e-16}{\coulomb}$.
\item C: poids confondu avec la masse ($|q|=m/E$, facteur $g$ oublié) : $\qty{1.6e-19}{\coulomb}$.
\item D: $E=U\times d$ au lieu de $U/d$ : $|q|=mg/(U d)=\qty{4.0e-15}{\coulomb}$.
\end{itemize}
\vspace{8pt}\hrule\vspace{8pt}

% ================= Q10 =================
\noindent\textbf{PHYS-F09-Q10 --- Quantification de la charge : nombre d'électrons}\par
Un corps porte une charge électrique négative de valeur absolue $|q|=\qty{4.8e-9}{\coulomb}$. Sachant que $e=\qty{1.6e-19}{\coulomb}$, le nombre d'électrons en excès portés par ce corps est :
\begin{enumerate}[label=\Alph*)]
\item $\num{3.0e-28}$
\item $\num{3.0e-10}$
\item $\num{3.0e10}$
\item $\num{3.0e28}$
\end{enumerate}
\textit{Bonne réponse : C}\par
La charge est un multiple entier de la charge élémentaire : $|q|=N e$, donc $N=\dfrac{|q|}{e}=\dfrac{\num{4.8e-9}}{\num{1.6e-19}}=\num{3.0e10}$ électrons. (Mantisse : $4{,}8/1{,}6=3{,}0$ ; exposant : $-9-(-19)=+10$.)\par
\textit{Pièges :}
\begin{itemize}
\item A: charges multipliées au lieu d'être divisées ($|q|\times e$), exposants additionnés.
\item B: division inversée $e/|q|$ (exposant $-19-(-9)=-10$).
\item D: exposants additionnés en valeur absolue lors de la division ($9+19=28$).
\end{itemize}
\vspace{8pt}\hrule\vspace{8pt}

\end{document}
